\chapter{Conversations That Do Not End}

\section*{7.1}
Yes: $\mu t.!\Nat.t$ is guarded. No: $\mu t.t$ is unguarded.

\section*{7.2}
$?\Nat.!\Bool.(\mu t.?\Nat.!\Bool.t)$.

\section*{7.3}
Unfolding rewrites a type expression to expose its next constructor. It does not schedule or execute a process.

\section*{7.4}
$\mu t.?Req.!Resp.t$.

\section*{7.5}
$\mu t.t$ is rejected because the recursive occurrence is reached without crossing a communication constructor.

\section*{7.6}
First unfolding: $!Ping.?Pong.R$ where $R=\mu t.!Ping.?Pong.t$. Second unfolding exposes the same $!Ping.?Pong$ prefix again after the recursive occurrence.

\section*{7.7}
A server can satisfy every finite prefix of an intended request/reply discipline while intentionally remaining available for further requests.

\section*{7.8}
T4.10 states only that unfolding preserves well-formed syntax. It has no scheduler, response-time, or eventuality premise.

\section*{7.9}
For $\mu t.!A.S$, guardedness places recursive occurrences in $S$ beneath the send constructor. Substitution re-inserts the well-formed recursive type in those positions; send formation is preserved.

\section*{7.10}
T4.10 legitimizes the unfolding that exposes a finite outer prefix; T4.4 then preserves the principal communication step through that prefix.

\section*{7.11}
One possible contract is: under weak fairness, every continuously enabled request-handling action is eventually scheduled. This is an added temporal assumption, not guardedness.

\section*{7.12}
Infinite correct trace: repeated `recv Req; send Resp`. Finite stuck trace: a guarded peer waiting on a message that its environment never sends. Guarded syntax permits both without extra liveness assumptions.
